\documentclass[12pt,a4paper]{article}
\usepackage[margin=2cm]{geometry}
\usepackage[spanish]{babel}
\usepackage{graphicx, multicol, latexsym, amsmath, amssymb}
\usepackage{array}
\usepackage{tabularx}
\usepackage{setspace}
\usepackage{titlesec}
\usepackage{subcaption}
\usepackage{float}
\usepackage{url}
\usepackage{hyperref}

% Ajuste de espacio entre secciones
\titlespacing*{\section}{0pt}{0.5em}{0.5em}

\renewcommand{\arraystretch}{1.3}

\begin{document}

% ---------------------------
% Encabezado con datos
% ---------------------------
\noindent
\begin{tabularx}{\textwidth}{|X|X|}
\hline
\textbf{Carné:} 2020045294 & \textbf{Nombre:} Justin Jaffeth Corea Masís \\
\hline
\textbf{Correo:} coreajustin288@estudiantec.cr & \textbf{Semana:} 6 del 10/09/2025 al 17/09/2025 \\
\hline
\end{tabularx}

\vspace{0.5cm}

\section{Resultados}
\begin{figure}[H]
	\centering
	\begin{subfigure}[b]{0.45\textwidth}
		\includegraphics[width=\textwidth]{/home/akai/Downloads/W&B Chart 9_16_2025, 9_20_52 PM.png}
	\end{subfigure}
	\hfill
	\begin{subfigure}[b]{0.45\textwidth}
		\includegraphics[width=\textwidth]{/home/akai/Downloads/W&B Chart 9_16_2025, 9_20_40 PM.png}
	\end{subfigure}

	\begin{subfigure}[b]{0.45\textwidth}
		\includegraphics[width=\textwidth]{/home/akai/Downloads/W&B Chart 9_16_2025, 9_20_07 PM.png}
	\end{subfigure}
	\hfill
	\begin{subfigure}[b]{0.45\textwidth}
		\includegraphics[width=\textwidth]{/home/akai/Downloads/W&B Chart 9_16_2025, 9_20_22 PM.png}
	\end{subfigure}
	
	\caption{Resultados de entrenamiento de los modelos}
\end{figure}

El mejor modelo tiene casi 91.5\% de precisión en el reconocimiento de placas en validación.
No hay sobreajuste, puesto que el error de validación no aumenta al disminuir el error de entrenamiento.
Estoy tratado de mejorar el modelo para tener más precisión.

De momento puede reconocer bastante bien placas de motos (una y dos filas), vehículos livianos, vehículos del TEC,
vehículos particulares y Taxis de San José. Otros casos de vehículos no los reconoce bien debido a
la falta de datos para entrenar ni validar. Me centré en esos casos que creo son los más probables
a ver en los estacionamientos del Campus en San José y San Carlos.

\section{Ejemplos de funcionamiento}
\begin{figure}[H]
	\begin{subfigure}[b]{0.45\textwidth}
		\includegraphics[width=\textwidth]{/home/akai/Pictures/app-ejemplo-1.png}
	\end{subfigure}
	\hfill
	\begin{subfigure}[b]{0.45\textwidth}
		\includegraphics[width=\textwidth]{/home/akai/Pictures/app-ejemplo-2.png}
	\end{subfigure}
	
	\begin{subfigure}[b]{0.45\textwidth}
		\includegraphics[width=\textwidth]{/home/akai/Pictures/app-ejemplo-3.png}
	\end{subfigure}
	\hfill
	\begin{subfigure}[b]{0.45\textwidth}
		\includegraphics[width=\textwidth]{/home/akai/Pictures/app-ejemplo-4.png}
	\end{subfigure}
	\caption{Ejemplos de detección y reconocimiento de placas vehiculares}
\end{figure}

\end{document}
